% !TeX root = ../sections/02_exercises.tex
\subsection{2-A-1.}
\begin{exercise}
	\(A,B\subset\mathbb{R}\) は下に有界であるとする。
	以下の問いに答えよ。
	
	\begin{enumerate}
		\item
		\(i=\inf A\) であるための必要十分条件は，次の二条件：
		\begin{enumerate}
			\item
			任意の \(a\in A\) に対して
			\[
			a\geq i
			\]
			である。
			
			\item
			任意の \(\varepsilon>0\) に対して，
			ある \(a_\varepsilon\in A\) が存在して
			\[
			i+\varepsilon>a_\varepsilon
			\]
			である。
		\end{enumerate}
		が成立することであることを証明せよ。
		
		\item
		\(A\subset B\) ならば
		\[
		\inf B\leq \inf A
		\]
		であることを証明せよ。
	\end{enumerate}
\end{exercise}

\begin{background}
	この問題の中心は，\(\inf A\) の意味を
	「\(A\) の下界のうち最大のもの」として理解することである。
	
	まず，実数 \(l\) が \(A\) の下界であるとは，
	任意の \(a\in A\) に対して
	\[
	l\leq a
	\]
	が成り立つことをいう。
	つまり，下界とは \(A\) のすべての元より下にある数である。
	
	したがって，
	\[
	i=\inf A
	\]
	とは，次の二つを同時に満たすということである。
	
	\begin{enumerate}
		\item
		\(i\) 自身は \(A\) の下界である。
		
		\item
		\(i\) より少しでも大きい数は，もはや \(A\) の下界ではない。
	\end{enumerate}
	
	ここで重要なのは，二つ目の条件である。
	\(i+\varepsilon\) が \(A\) の下界ではないとは，
	ある \(a_\varepsilon\in A\) が存在して
	\[
	a_\varepsilon<i+\varepsilon
	\]
	となることを意味する。
	
	つまり，
	\[
	i=\inf A
	\]
	であることは，
	\[
	i\leq a \quad (a\in A)
	\]
	によって「\(i\) が下界である」ことを保証し，
	さらに
	\[
	\forall \varepsilon>0,\quad
	\exists a_\varepsilon\in A
	\quad\text{s.t.}\quad
	a_\varepsilon<i+\varepsilon
	\]
	によって「\(i\) より大きい数は下界になれない」
	ことを保証している。
	
	また，\(A\subset B\) のとき，\(B\) のすべての元以下である数は，
	特に \(A\) のすべての元以下でもある。
	したがって，
	\[
	\{\text{\(B\) の下界}\}
	\subset
	\{\text{\(A\) の下界}\}
	\]
	となる。
	下界の集合が小さくなると，その最大元である下限は小さくなる。
	これが
	\[
	\inf B\leq \inf A
	\]
	の発想である。
\end{background}

\begin{strategy}
	この問題では，下限を直接計算するのではなく，
	\[
	\inf A=\max\{\text{\(A\) の下界}\}
	\]
	という定義を，\(\varepsilon\) を使った形に言い換える。
	
	\begin{enumerate}
		\item[(1)]
		まず，
		\[
		i=\inf A
		\]
		を仮定する。
		このとき \(i\) は \(A\) の下界なので，
		任意の \(a\in A\) に対して
		\[
		i\leq a
		\]
		が成り立つ。
		したがって条件 (i) が従う。
		
		次に，任意に \(\varepsilon>0\) を取る。
		ここで見るべき数は
		\[
		i+\varepsilon
		\]
		である。
		これは \(i\) より大きい数である。
		
		もし \(i+\varepsilon\) が \(A\) の下界であれば，
		\(i\) より大きい下界が存在することになる。
		これは \(i\) が下界の最大元であることに反する。
		
		したがって，\(i+\varepsilon\) は \(A\) の下界ではない。
		下界ではないということは，ある \(a_\varepsilon\in A\) が存在して
		\[
		a_\varepsilon<i+\varepsilon
		\]
		となることである。
		よって条件 (ii) が従う。
		
		逆に，条件 (i), (ii) を仮定する。
		条件 (i) より，\(i\) は \(A\) の下界である。
		あとは，\(i\) が下界の中で最大であることを示せばよい。
		
		そのために，任意に
		\[
		j>i
		\]
		を取る。
		この \(j\) が \(A\) の下界ではないことを示す。
		
		ここで
		\[
		\varepsilon=j-i>0
		\]
		とおく。
		条件 (ii) より，ある \(a_\varepsilon\in A\) が存在して
		\[
		a_\varepsilon<i+\varepsilon=j
		\]
		となる。
		つまり，
		\[
		a_\varepsilon<j
		\]
		である。
		
		したがって，\(j\) は \(A\) のすべての元以下ではない。
		よって \(j\) は \(A\) の下界ではない。
		
		以上より，\(i\) は \(A\) の下界であり，
		\(i\) より大きい数はどれも下界ではない。
		したがって
		\[
		i=\inf A
		\]
		である。
		
		\item[(2)]
		\(A\subset B\) とする。
		このとき，\(B\) の下界を一つ取る。
		それを \(l\) とする。
		
		\(l\) が \(B\) の下界であるとは，任意の \(b\in B\) に対して
		\[
		l\leq b
		\]
		が成り立つことである。
		
		いま \(A\subset B\) だから，任意の \(a\in A\) は \(B\) の元でもある。
		したがって
		\[
		l\leq a
		\]
		が任意の \(a\in A\) に対して成り立つ。
		つまり，\(l\) は \(A\) の下界でもある。
		
		よって，
		\[
		\{\text{\(B\) の下界}\}
		\subset
		\{\text{\(A\) の下界}\}
		\]
		である。
		
		ここで，下限は下界全体の最大元である。
		\(B\) の下界全体は，\(A\) の下界全体に含まれているので，
		その最大元は \(A\) の下界全体の最大元を超えられない。
		したがって
		\[
		\inf B\leq \inf A
		\]
		が成り立つ。
	\end{enumerate}
\end{strategy}